\section{2/25}
Most applications are reactive and they can be described as \textbf{state machines}. A state machine is a set of states and transitions between them. A common example is a vending machine.
\\
Some state machines have an unbounded amount of states, these can be represented by \textbf{process algebras}. These are important for reasoning about systems that change.
\begin{lstlisting}
module type StateMachine = sig
  type state
  type input
  val init : state
  val step : state -> input -> state
end
(* 
INIT := MACHINE(0.00, 10)
  MACHINE(money, items) :=
    insert25c.MACHINE(money + 0.25, items)
    + insert50c.MACHINE(money + 0.50, items)
    + vend.MACHINE(money - 1.00, items - 1) if money >= 1.00 and items >= 1
*)

module MoneyItemVendingMachine (* StateMachine *) = struct
  type state = { money: float; items: int }
  type input = Insert25c | Insert50c | Vend

  let init = { money= 0.0; items= 10}
  
  let step state input = 
    match input with
    | Insert25c -> { state with money = state.money +. 0.25 }
    | Insert50c -> { state with money = state.money +. 0.50 }
    | Vend when state.money >= 1.00 && state.items >= 1 -> 
      { money = state.money -. 1.00; items = state.items - 1}
    | Vend when state.money < 1.00 -> failwith "not enough money"
    | Vend when state.items < 1 -> failwith "out of stock"
    | Vend -> failwith "invalid state"
end
\end{lstlisting}
In order to test functions that change sequentially, we need a trace. Generate a sequence of inputs and consider properties about the state over the sequence.
\\\\
We can write a functor that takes in a state machine and turns it into a machine trace. It takes in a list of inputs. The entire struct is an option as if an error is thrown we return \textbf{None}.